\scanpage{13}
Sous ces conditions, la factorisation (i) est unique à isom près, et $S'$ est de \deleted{t.f.} \inserted{prés. finie} sur $S$.

$(i)\Rightarrow(i\ \mathrm{bis})\Rightarrow(ii)$ triv. De $=$ l'unicité de la factorisation \inserted{(i)}, car $X\longrightarrow S'$ \unsure{surj. quac.}, donc couvrant pour fpqc donc $S'=X/R$, où $R=X\times_S X$.

Reste à prouver $(ii)\Rightarrow$ \inserted{le faisceau} $X/R$ est représentable par un $S'$ de \inserted{loc.} type fini sur $S$. \deleted{[N.B. alors \unsure{cela} \unsure{quasi compact} sur $S$ car $X\longrightarrow S'$ est}

[N.B. Alors $X\longrightarrow S'$ est couvrant pour fpqc, comme \unsure{après le chgt|on voit du. chgt} de base $X-S'$ \unsure{il|et} donne $X\times_S X=R\longrightarrow X$ qui est fid plat par hypothèse ; il est fid plat \unsure{par. cqt} lui-$=$, ce qui implique \inserted{aussi} que $S'$ est \inserted{\unsure{par. cqt loc.}} de t.f. \deleted{\unsure{donc quasi}\ldots})].

Or \deleted{\unsure{comme} $X/R\longrightarrow S$} le foncteur représenté par $X/R\unsure{\hookrightarrow|\leftarrow}S$ est ici le suivant :
\[
F(\unsure{T})=\begin{cases}\varnothing&\text{si }X_T\text{ non fid plat sur }T\ \inserted{[\Leftrightarrow X_T\longrightarrow T\text{ couvrant fpqc}]}\\\{\varnothing\}&\text{sinon}\end{cases}
\]
D'ailleurs on est ramené au cas $S$ affine, \uuline{noethérien}.
\deleted{et} Alors on applique le th., qui nous dit qu'il suffit de vérifier \deleted{\unsure{que}} $(L)$. Or sous la condition de cette condition, $X_T\longrightarrow T$ est couvrant [car il contient un pt \inserted{$x$} de la fibre spéciale, et $\operatorname{Sp}\underline{O}_{X,x}-\operatorname{Sp}\underline{O}_{T,t_0}$ est fid plat] donc on \unsure{gagne} !
\mspage{12}
